\section{Introduction}

The rapid proliferation of the Web3 ecosystem and cryptocurrencies has ushered in a new era of digital markets, bringing unprecedented opportunities for innovation, efficiency, and financial inclusion \cite{bonneau2015sokresearchperspectives,voshmgir2020tokeneconomyhow,abdelhamid2025reviewblockchaintechnology}. Yet the very attributes that empower these technologies, including decentralization, pseudonymity, and global accessibility, simultaneously pose profound challenges to Anti-Money Laundering (AML) and Countering the Financing of Terrorism (CFT) compliance \cite{force2021updatedguidanceriskbased,financialactiontaskforcefatf2012internationalstandardscombating,huang2020antimoneylaunderingblockchain}. The inherent characteristics of blockchain systems have exposed the limitations of traditional compliance frameworks, which were primarily designed for centralized financial infrastructures, rendering them ill-suited for the complexities of the digital asset landscape \cite{weber2019antimoneylaunderingbitcoin,pocher2023detectinganomalouscryptocurrency,bankforinternationalsettlementsbis2022futuremonetarysystem,aldasoroapproachantimoneylaundering}. In response, a new generation of regulatory technology (RegTech) solutions has emerged, purpose-built to address the distinct compliance demands of the Web3 environment \cite{arner2016fintechregtechreconceptualization,charoenwong2024regtechtechnologydrivencompliance,prosper2025blockchainbasedsolutionsantimoney}.

Despite significant progress in blockchain compliance capabilities, formidable challenges remain. The cryptocurrency compliance landscape exists in a constant state of flux, with illicit actors and compliance providers locked in perpetual arms race dynamics \cite{albrecht2019usecryptocurrenciesmoney,guidara2022cryptocurrencymoneylaundering,prosper2025blockchainbasedsolutionsantimoney}. Novel evasion tactics emerge rapidly as adversaries adapt to new detection methods, requiring RegTech solutions to maintain continuous innovation cycles \cite{unitednationsofficeondrugsandcrimemoneylaunderingcryptocurrencies}. Compliance providers respond by developing advanced analytical capabilities leveraging graph neural networks, AI-powered behavioral analytics, and sophisticated forensic techniques \cite{azad2025machinelearningblockchain,rodriguezvalencia2025systematicreviewartificial,zebari2025comprehensivereviewintegrating}. This adversarial co-evolution shows no signs of abating, necessitating adaptive RegTech solutions that can evolve alongside emerging threats \cite{verma2024impactcryptocurrencymoney,tiwari2025understandingevolvinglandscape}.

The absence of universally accepted standards poses significant challenges across multiple compliance domains. Risk scoring methodologies for addresses, transactions, and entities lack standardization, hindering Virtual Asset Service Providers' ability to benchmark compliance efforts and limiting regulators' capacity to compare different approaches objectively. The FATF Travel Rule mandates information exchange between VASPs, yet the absence of a single universally adopted protocol for secure and interoperable data sharing has resulted in fragmented solutions and operational complexities \cite{force2021updatedguidanceriskbased,pocher2022privacytransparencycbdcs}. Rapid blockchain innovation continually introduces new protocols, token standards, Layer-2 solutions, and decentralized applications, creating moving targets for compliance tools \cite{mohammedabdul2024navigatingblockchainstwin,xu2024exploringblockchaintechnology,bankforinternationalsettlementsbis2022futuremonetarysystem}. Each technological advancement introduces unique risks requiring constant analytical adaptation. The growing adoption of DeFi protocols and self-custody wallets further complicates compliance by limiting VASP visibility into customer activities, creating blind spots for AML/CFT monitoring and Travel Rule implementation when funds move off-platform \cite{harvey2021defifuturefinance,salami2021challengesapproachesregulating}.

\subsection{Research Questions and Contributions}

This paper systematizes knowledge about Web3 RegTech solutions and the regulatory compliance lifecycle. Information about these systems has been scattered across industry whitepapers, academic publications, regulatory guidance documents, and practitioner reports. We provide foundational background on compliance regulation, analyze current solutions across industry and academia, and identify critical gaps and research directions. We organize our contributions around three research questions.

\noindent\textbf{RQ1: What are the fundamental architectural differences between Web3 and traditional RegTech solutions, and what novel technical capabilities do blockchain-native approaches enable for VASP AML/CFT compliance?}
We systematically compare centralized and blockchain-native compliance architectures across seven dimensions: data accessibility, verification mechanisms, identity models, monitoring scope, analytical approaches, enforcement models, and transparency guarantees \cite{arner2016fintechregtechreconceptualization,zetsche2017fintechtechfinregulatory}. Our analysis demonstrates that Web3 RegTech constitutes a paradigmatic shift rather than incremental evolution, enabling comprehensive transaction graph analysis \cite{meiklejohn2013fistfulbitcoinscharacterizing,reid2013analysisanonymitybitcoin}, real-time risk assessment \cite{moser2014riskscoringbitcoin,weber2019antimoneylaunderingbitcoin}, cross-chain analytics \cite{yousaf2019tracingtransactionscryptocurrency,zamyatin2021sokcommunicationdistributed}, smart contract interaction analysis \cite{chen2020understandingethereumgraph,qin2021attackingdefiecosystem}, and privacy-preserving verification \cite{ben-sasson2017scalablezeroknowledge,zhang2024blockchainzeroknowledge}. We contextualize these architectural differences with recent advances in consensus diagnostics, hybrid security, and AI-optimized blockchain controls that increasingly intersect with compliance requirements \cite{ahn2024blockchainconsensusmechanisms,cai2025blockchainconsensusalgorithm,wei2024enhancedconsensusalgorithm,venkatesan2024blockchainsecurityenhancement,yuan2025aidrivenoptimizationblockchain}. Through examination of 41 operational RegTech platforms (34 Web3-native and 7 traditional with blockchain capabilities), we document rapid market growth with numerous Web3-native solutions emerging since 2023. \autoref{sec:systematization} and \autoref{sec:ecosystem} address this research question.

\noindent\textbf{RQ2: What is the current landscape of Web3 RegTech across industry and academia, what are their technical approaches, and what gaps exist between proposed solutions and practical VASP compliance needs?}
We analyze 41 commercial Web3 RegTech platforms and 28 academic research prototypes, categorizing them into five functional domains: transaction monitoring \cite{weber2019antimoneylaunderingbitcoin,petterssonruiz2022combatingmoneylaundering}, risk scoring \cite{moser2014riskscoringbitcoin,monamo2016unsupervisedlearningrobust}, blockchain forensics \cite{meiklejohn2013fistfulbitcoinscharacterizing,yousaf2019tracingtransactionscryptocurrency}, compliance automation \cite{charoenwong2024regtechtechnologydrivencompliance}, and privacy-preserving verification \cite{ben-sasson2017scalablezeroknowledge,yang2024researchidentitydata}. We examine technical building blocks including graph neural networks \cite{wu2023tracerscalablegraphbased,azad2025machinelearningblockchain,asiri2025graphconvolutionnetwork,zhang2024graphnetworkmodels}, clustering algorithms \cite{reid2013analysisanonymitybitcoin,ron2013quantitativeanalysisfull}, machine learning models \cite{pourhabibi2020frauddetectionsystematic,rodriguezvalencia2025systematicreviewartificial}, heuristic detection \cite{difrancescomaesa2018datadrivenanalysisbitcoin}, and cryptographic protocols \cite{zhang2024blockchainzeroknowledge,camenisch2001efficientsystemnontransferable}. Our gap analysis identifies eight critical disconnects between academic innovation and industry adoption: technology readiness limitations, regulatory alignment challenges, data availability constraints, adversarial robustness gaps, cost-benefit misalignments, standardization deficiencies, operational workflow integration failures, and validation benchmarking inadequacies. We document six persistent challenges limiting current capabilities across cross-chain tracking \cite{zamyatin2021sokcommunicationdistributed,lin2024crosschainabnormaltransaction}, DeFi interaction analysis \cite{qin2021attackingdefiecosystem,sun2025ethereumfrauddetection}, privacy protocol analysis \cite{kappos2018empiricalanalysisanonymity,kumar2017traceabilityanalysismoneros}, real-time scalability, attribution database maintenance, and false positive management \cite{alotibi2022moneylaunderingdetection}. \autoref{sec:ecosystem} addresses this research question.

\noindent\textbf{RQ3: Based on identified architectural differences and ecosystem gaps, what are the critical research directions and best practices to advance Web3 RegTech effectiveness?}
Building on comprehensive analysis of current capabilities and limitations, we identify critical research directions encompassing verifiable compliance proofs, semantic transaction understanding, behavior-based risk assessment, privacy-preserving collaborative analytics, intent-aware monitoring systems, cross-chain attribution techniques, adversarial robustness enhancements, standardized evaluation frameworks, and regulatory-compatible privacy preservation. We present architectural best practices for scalable monitoring, privacy-computation trade-offs, and intent interpretation frameworks that developers, protocol designers, and researchers can leverage to build solutions enhancing compliance effectiveness without compromising Web3's core principles of decentralization, transparency, and user sovereignty. \autoref{sec:discussion} addresses this research question.

\subsection{Research Scope and Methodology}
This systematization draws on analysis of 41 operational RegTech platforms (34 Web3-native, 7 traditional with blockchain capabilities) and 28 academic research prototypes from systematic literature review (2015-2025). Following established methodology for technology ecosystem analysis \cite{arner2016fintechregtechreconceptualization,bonneau2015sokresearchperspectives}, we synthesize commercial platform capabilities through technical documentation, whitepapers, and public capability descriptions, complemented by academic literature providing independent validation. This mixed-methods approach enables comprehensive domain coverage while acknowledging the inherent challenges of analyzing rapidly evolving technology ecosystems where platform capabilities continuously advance.

% - **RQ1**: 聚焦Web3与传统RegTech的架构差异和技术能力对比
% - **RQ2**: 重构为对当前RegTech生态（学术+产业）的系统性分析
% - **RQ3**: 整合为基于前面RQ的最佳实践和研究方向


% \noindent\textbf{A taxonomy of RegTech countermeasures.}
